\newpage
\section{Conclusion}\label{s:conclusion}

This review highlighted the adaptability and efficacy of angiogenic therapies in AMD prevention. \TreatAndExtend regimens have emerged as an effective alternative to traditional approaches, offering comparable visual outcomes with a reduced treatment burden. Combination therapies show potential for enhancing outcomes, though further research is needed.

Safety profiles of treatments like ranibizumab and bevacizumab were generally favorable, supporting their use in diverse clinical settings. These findings provide clinicians with effective, safe options for AMD treatment, allowing for tailored approaches that could improve patient adherence and satisfaction.

Our review contributes to the ongoing discussion on the optimal use of angiogenic therapies in AMD prevention. It underscores the need for continuous research to refine treatment protocols, accommodate patient variability, and explore new therapeutic modalities. Informed by this systematic review, clinicians can navigate the complexities of AMD care, aiming to improve visual outcomes and enhance patients' quality of life.